\section{Introduction}\label{sec:intro}
\PARstart{A}{utonomous} vehicles and mobile robots carry several sensors because each one fails in different conditions. A camera loses contrast at night and in fog, a depth estimate develops holes on reflective surfaces, an event camera goes silent on static scenes, and LiDAR returns thin out in rain and with distance. Multimodal semantic segmentation promises to cover these failures by fusing the sensors, so that the scene is still understood when one of them degrades.

Research on this problem has split into two tracks. The first track maximizes accuracy when all sensors are present~\cite{cmnext2023,cmx2023,stitchfusion2024,geminifusion2024,cafuser2025,dgfusion2026}. The second track measures and improves robustness when sensors are missing or noisy~\cite{magic2024,anyseg2024,rmmss2025,mmssbench2025,fuseadapt2026}. The two tracks rarely evaluate the same models under the same protocol. As a result, it is unclear whether the most accurate fusion models are also robust, and if they are not, what makes them fragile. A recent robustness benchmark~\cite{mmssbench2025} reports that several fusion models lose much of their accuracy when a sensor is removed, and it explains the differences by fusion design, contrasting RGB-centric, equal-contribution, and adaptive-selection architectures. Whether the design is in fact the cause has not been tested by holding the design fixed and changing only the training.

We study this question on three benchmarks with complementary sensor suites: DELIVER~\cite{cmnext2023} (RGB, depth, event, LiDAR), MUSES~\cite{muses2024} (RGB, event, LiDAR, radar), and MCubeS~\cite{mcubes2022} (RGB, polarization, near-infrared). We evaluate a leading accuracy-oriented model, DGFusion~\cite{dgfusion2026}, and our own fusion model under one protocol, and find that both rely heavily on depth on DELIVER: removing depth costs 23 to 39 points of mIoU, more in the dark than in daylight. DGFusion uses depth supervision, which might explain its dependence, but our model has no such supervision and shows the same pattern. When either model is trained with synthetically degraded inputs in addition to clean ones, the dependence largely disappears. The over-reliance is therefore learned from clean training data rather than built into a specific architecture.

Building on this finding, we present \ours{}, a fusion model on a frozen DINOv3 backbone with sensor-specific adapters, symmetric cross-modal attention, and a high-resolution decoding path, trained with a clean pass and a degraded pass that is distilled from a frozen clean teacher. The training scheme needs no condition or degradation labels and adds no inference cost. On DELIVER, \ours{} reaches 57.74 mIoU on the test set (mean of three seeds), 1.0 point above DGFusion, while reducing the depth over-reliance to a third or less, and it scores 57.85 under the standard missing-modality protocol on the validation set, 7.9 points above the best published result \todo{add the MUSES and MCubeS results once measured}.
% Sources: headline.yaml deliver.ours[0] (57.74, gap +1.03 vs DGFusion 56.71); tasks/DRN-261019-03/analysis.md (val EMM15 57.85 vs RMMSS 49.93, +7.92).

Our contributions are as follows.
\begin{itemize}
\item We show, with controlled comparisons on the same backbone and with re-trained baselines evaluated under one protocol, that the sensor over-reliance of accurate fusion models on DELIVER comes from clean training rather than from their architecture or supervision.
\item We combine a frozen foundation backbone with a two-pass training scheme and a high-resolution path into a single model that improves clean accuracy and robustness together, instead of trading one for the other.
\item We analyse per-sensor dependence across conditions, missing, partial, and noisy sensors, thin objects, and failure cases, and we report where the model still fails. \todo{state the code and checkpoint release}
\end{itemize}
